% ============================================================
%  LaTeX Cheat-Sheet  --  built from the CSE200 Online-1 files
%  (Q1 Image Synthesis, Q2 Intro, Q3 Fireflies, Q4 Laplace, Q5 Solar System)
%  Compile:  pdflatex LaTeX_Cheatsheet.tex   (run twice for the \ref numbers)
%  Every snippet below is real syntax pulled from your own documents.
% ============================================================
\documentclass[10pt]{article}

\usepackage[a4paper,landscape,margin=1cm]{geometry}
\usepackage{multicol}
\usepackage[T1]{fontenc}
\usepackage[utf8]{inputenc}
\usepackage{amsmath}
\usepackage{amssymb}
\usepackage{xcolor}
\usepackage{listings}
\usepackage{parskip}

\definecolor{headcol}{RGB}{40,80,170}
\definecolor{codebg}{RGB}{245,246,250}
\definecolor{cmdcol}{RGB}{170,40,40}

% ---- style for the LaTeX-source snippets --------------------
\lstset{
    language=[LaTeX]TeX,
    basicstyle=\ttfamily\footnotesize,
    keywordstyle=\color{cmdcol},
    commentstyle=\color{gray}\itshape,
    backgroundcolor=\color{codebg},
    breaklines=true,
    columns=fullflexible,
    keepspaces=true,
    xleftmargin=3pt, xrightmargin=3pt,
    aboveskip=3pt, belowskip=3pt,
    frame=none,
    showstringspaces=false,
}

% ---- compact coloured section heading -----------------------
\newcommand{\hd}[1]{%
    \par\medskip
    {\color{headcol}\normalsize\bfseries #1}\par
    {\color{headcol}\hrule height 1pt}\smallskip}

\setlength{\columnsep}{18pt}
\pagestyle{empty}

\begin{document}
\begin{center}
    {\color{headcol}\LARGE\bfseries LaTeX Cheat-Sheet}\quad
    {\small — grounded in your CSE200 Online-1 documents}
\end{center}

\begin{multicols}{2}

% ============================================================
\hd{1. Document Skeleton}
The frame every one of your files shares.
\begin{lstlisting}
\documentclass[11pt]{article}
\usepackage[T1]{fontenc}
\usepackage[utf8]{inputenc}
\usepackage[margin=1in]{geometry}
\usepackage{amsmath, amssymb, mathtools}
\usepackage{graphicx, xcolor, enumitem}

\title{Document Title}
\author{Name\\ \small Affiliation}
\date{}                 % empty -> no date
\begin{document}
\maketitle
   ... content ...
\end{document}
\end{lstlisting}
\texttt{\%} starts a comment. Blank line = new paragraph.

% ============================================================
\hd{2. Sectioning}
\begin{lstlisting}
\section{The Puzzle in the Forest}
\subsection{A Firefly as a Clock}
\subsection*{Circuit Model}   % * = unnumbered
\section*{Conclusion}         % no number, no TOC
\tableofcontents \newpage     % needs 2nd compile
\end{lstlisting}
Rename a heading (Q5): \verb|\renewcommand{\contentsname}{Contents}|

% ============================================================
\hd{3. Text \& Fonts}
\begin{lstlisting}
\textbf{bold}   \textit{italic}   \emph{emphasis}
\underline{underlined}   \texttt{monospace}
\textsc{Small Caps}
{\Large big} {\small small} {\huge huge}
\textcolor{blue}{coloured words}
super\textsuperscript{1}   x\textsubscript{0}
\LaTeX   \qquad   ~ (non-break space)
--- (em dash)   -- (en/range)   ` ' quotes
\end{lstlisting}
Size steps: \verb|\tiny \scriptsize \footnotesize \small|
\verb|\normalsize \large \Large \LARGE \huge \Huge|

% ============================================================
\hd{4. Lists  (enumitem)}
\begin{lstlisting}
\begin{itemize}
    \item First point
    \item[\textbf{Note}] custom label
\end{itemize}

\begin{enumerate}[label=\Roman*.]   % I. II. III.
    \item Drift
\end{enumerate}
% other labels: (\Roman*)  (\alph*)  \arabic*)

\begin{itemize}[label=$\diamond$]   % custom bullet
    \item ...
\end{itemize}

\begin{description}[style=nextline]
    \item[Coupling] explanation on next line
\end{description}
\end{lstlisting}
Handy keys: \verb|nosep|, \verb|leftmargin=2em|, \verb|label=$\oplus$|.
Lists nest freely (see Q1/Q2).

% ============================================================
\hd{5. Inline \& Display Math}
\begin{lstlisting}
Inline: $\theta_i(t) \in [0,2\pi)$, $x_i$, $z^2$.

Unnumbered display:
\[  \Delta = \theta_2 - \theta_1  \]

Numbered, with a label:
\begin{equation}\label{eq:kuramoto}
    T_i = \frac{2\pi}{\omega_i}
\end{equation}
\end{lstlisting}

% ============================================================
\hd{6. Multi-line Math}
\begin{lstlisting}
% aligned at &, each line numbered
\begin{align}
    c^2 &= a^2 - b^2 \\
    e   &= \frac{c}{a} \notag   % \notag = no number
\end{align}
\begin{align*} ... \end{align*}  % none numbered

% one number for a multi-line derivation
\begin{equation}
  \begin{split}
    \mathcal{L}\{1\} &= \int_0^\infty e^{-st}\,dt \\
                     &= \frac{1}{s}
  \end{split}
\end{equation}

% piecewise
f(x)=\begin{cases}
    x^2/\sigma^2 & \text{if } x \ge 0 \\
    |x|/\sigma   & \text{if } x < 0
\end{cases}
\end{lstlisting}

% ============================================================
\hd{7. Math Symbols You Use}
\begin{lstlisting}
Greek:  \theta \omega \sigma \mu \pi \alpha
        \beta \gamma \Delta \Lambda \varpi \varepsilon
Ops:    \sum_{j=1}^{N}  \int_0^\infty  \prod
        \frac{a}{b}  \dfrac{a}{b}  \sqrt{x}
        \frac{\mathrm{d}\theta}{\mathrm{d}t}
Rel:    \le \ge \approx \propto \equiv \in \pmod
        \to \Rightarrow \implies \xrightarrow{L}
Deco:   \bar x \hat x \vec x \dot x  \bmod
        \boxed{|\Delta\omega| \le K}
Cal:    \mathcal{L}  \mathbb{R}  M_\odot  \Gamma^\mu_{\nu}
\end{lstlisting}
\verb|\,| thin space, \verb|\quad|/\verb|\qquad| wide space,
\verb|\text{...}| for words inside math.

% ============================================================
\hd{8. Tables}
\begin{lstlisting}
\begin{table}[ht]
  \centering
  \caption{Caption above or below}
  \label{tab:x}
  \begin{tabular}{|l|c|c|}      % l/c/r + | rules
    \hline
    \multirow{2}{*}{Planet}
      & \multicolumn{2}{c|}{Orbit} \\
    \cline{2-3}
      & $a$ & $e$ \\
    \hline
    Jupiter & 5.204 & 0.049 \\
    \hline
  \end{tabular}
\end{table}
\end{lstlisting}
\textbf{booktabs} (Q3): \verb|\toprule \midrule \bottomrule| (no vertical bars).\\
\textbf{Coloured rows} (needs \verb|\usepackage[table]{xcolor}|):
\verb|\rowcolor{green!10}|\\
\textbf{Wrapping cell}: column type
\verb|>{\raggedright\arraybackslash}p{4.3cm}|

% ============================================================
\hd{9. Figures}
\begin{lstlisting}
\begin{figure}[ht]
    \centering
    \includegraphics[width=0.6\textwidth]{img.png}
    \caption{...}\label{fig:x}
\end{figure}
\end{lstlisting}
\textbf{Text wraps around} (wrapfig, Q3/Q4):
\begin{lstlisting}
\begin{wrapfigure}{r}{0.42\textwidth}
  \centering
  \includegraphics[width=0.40\textwidth]{f.jpg}
  \caption{...}\label{fig:f}
\end{wrapfigure}
\end{lstlisting}
\textbf{Side-by-side} (subcaption, Q1/Q5):
\begin{lstlisting}
\begin{subfigure}{0.3\textwidth}\centering
    \includegraphics[width=\textwidth]{a}
    \caption{Output A}
\end{subfigure}\hfill
\begin{subfigure}{0.3\textwidth} ... \end{subfigure}
\end{lstlisting}
\verb|\hfill| spreads them; blank line starts a new row.

% ============================================================
\hd{10. Cross-References}
\begin{lstlisting}
Give a \label{...} after the caption/equation,
then refer to it anywhere:
Table~\ref{tab:sine}      fig.~\ref{fig:f}
eq.~\eqref{eq:kuramoto}   % (adds parentheses)
Listing~\ref{lst:rc}
\end{lstlisting}
\verb|~| glues the word to the number so they never split across a line.
\emph{Always compile twice} to resolve numbers.

% ============================================================
\hd{11. Colour (xcolor)}
\begin{lstlisting}
\definecolor{titleblue}{RGB}{40,80,170}
\textcolor{titleblue}{words}
\textcolor{green!55!black}{Positive}   % mix
\rowcolor{red!10}                      % table row
red!70!black  gray!15  green!45!black  % blends
\end{lstlisting}

% ============================================================
\hd{12. Code Listings (Q4)}
\begin{lstlisting}
\usepackage{listings}
\lstdefinestyle{pystyle}{
    language=Python, basicstyle=\ttfamily\small,
    keywordstyle=\color{blue!70!black}\bfseries,
    numbers=left, frame=single, breaklines=true}

% pull code from an external file:
\lstinputlisting[style=pystyle,
    caption={...}, label=lst:rc]{rc_response.py}
\end{lstlisting}

% ============================================================
\hd{13. Theorems \& Proof (Q5)}
\begin{lstlisting}
\usepackage{amsthm}
\newtheorem{proposition}{Proposition}[section]

\begin{proposition}[Weak field]\label{prop:x}
    Statement of the result.
\end{proposition}
\begin{proof}
    Argument.   % gets the QED square
\end{proof}
\end{lstlisting}

% ============================================================
\hd{14. Algorithms (Q3)}
\begin{lstlisting}
\usepackage{algorithm, algpseudocode}
\begin{algorithm}[ht]
  \caption{Simulate the Kuramoto model}
  \begin{algorithmic}[1]
    \Require Phases, $K$, step $h$, steps $M$
    \For{$m \gets 1$ \textbf{to} $M$}
        \State $v_i \gets \omega_i + \frac{K}{N} q_i$
    \EndFor
    \State \Return $\theta_1,\dots,\theta_N$
  \end{algorithmic}
\end{algorithm}
\end{lstlisting}

% ============================================================
\hd{15. Links, Footnotes, Citations}
\begin{lstlisting}
\usepackage{hyperref}
\hypersetup{colorlinks=true, linkcolor=blue,
            urlcolor=blue, citecolor=black}
\href{https://...}{clickable text}
A remark here.\footnote{The note text.}

Reference~\cite{sarfati2021}.
\begin{thebibliography}{9}
  \bibitem{sarfati2021}
    Author. ``Title''. \emph{Journal} (2021).
\end{thebibliography}
\end{lstlisting}

% ============================================================
\hd{16. Spacing \& Boxes}
\begin{lstlisting}
\\[0.3em]  break + extra vertical space
\vspace{1em}  \bigskip \medskip \smallskip
\newpage  \par  \noindent  \centering
\hfill (push apart)   \quad \qquad (h-space)
\thispagestyle{empty}   % no page number
\end{lstlisting}

\end{multicols}
\end{document}
